\section{Introduction}\label{sec:introduction}
Finite-difference penalized least squares (PLS) is a classical approach to recovering a latent trend from noisy observations. Whittaker's quadratic-difference method \citep{whittaker1923}, the statistical formulation of \citet{guerrero2007}, and the broader Whittaker--Henderson tradition provide the foundation. A fixed nonzero mean of differences permits polynomial trends beyond the homogeneous special case. Guerrero's trace-based smoothness measure \citep{guerrero2008} and the comparison of classical optimality criteria by \citet{cortes2017} are likewise established starting points.

There are two different computational problems. First, at a given penalty $\lambda$, solve a structured positive-definite linear system to fit the trend. Second, when selecting a smoothing level, evaluate one or several criteria over many penalties. Cross-validation requires residuals and diagonal leverages, generalized cross-validation requires effective degrees of freedom, and a Gaussian likelihood needs quadratic forms and log determinants. Solving one banded system quickly does not establish the cost of evaluating an entire candidate grid. Conversely, paying for a dense spectral decomposition may be worthwhile when the result is reused.

The organizing variable in this paper is the user-facing normalized smoothness index $S\in[0,1]$. This reparameterization of the classical trace reduction is dimensionless and interpretable across orders: $S=0$ means interpolation; $S=1$ is the exactly constrained polynomial limit. A single SVD of the finite-difference operator supplies a monotone map between $S$ and the underlying penalty, fitted trends, and several candidate-scoring terms. The compact interval does not by itself guarantee globally optimal or faster numerical searches; its inversion has a cost, particularly near the endpoints.

The reference level $\mu$ is central to the formulation. The convenient assumption $\mu=0$ changes the limiting polynomial family and excludes prescribed nonzero $d$th-difference levels. While arbitrary fixed $\mu$ has already been treated by \citet{guerrero2007}, it is useful to have one numerically stable computational API handling both cases. In particular, explicitly forming a large degree-$d$ reference polynomial can lead to severe cancellation. Our shift-free singular-vector reconstruction avoids that intermediate.

The numerical and statistical foundations of efficient smoothing are extensive. \citet{weinert2007} exploits Whittaker--Henderson structure for efficient GCV, \citet{golub1979} develops classical spectral GCV methods, and \citet{biessy2026} supplies a modern Whittaker--Henderson statistical and computational treatment. We do not claim invention of diagonalization, the penalty, trace smoothness, cross-validation, Cholesky, or selected inversion. The contribution is a unified, reproducible \emph{workload-aware} treatment with explicit computational and numerical trade-offs, using the normalized smoothness coordinate throughout.

Specifically, the article develops (i) an arbitrary-order spectral formula supporting prescribed nonzero means without unstable polynomial shifts; (ii) the normalized index, its bijection and endpoints; (iii) the exact quantities required by conditional CV, GCV, information criteria and reduced Gaussian likelihood; (iv) time and storage bounds for complete spectral and banded pipelines across $M$ candidates, $R$ series and multiple criteria; and (v) a numerical-admissibility policy and open-source Python interface.

All principal theorems assume equally spaced observations, unweighted quadratic fidelity, integer $1\le d<n$, a fixed numeric $\mu$, and exact arithmetic. Data-derived reference means, time-series forecasting and fully measured hardware crossovers require additional statistical or empirical arguments, which are not silently incorporated into these claims.
